\section{Related Work}
\label{sec:related}

\myparagraph{Flaky tests} A flaky test produces different outcomes across runs of the same software version \cite{luo2014empirical, parry2021survey}. Developers spend considerable effort triaging failures that do not signal a real defect \cite{pinto2020vocabulary}, and repeated exposure erodes trust to the point that genuine failures are ignored \cite{gruber2022survey, rahman2018impact}; the cost to delivery pipelines is also documented \cite{herzig2015empirically, micco2017state}.

\myparagraph{Detection} Techniques are dynamic, hybrid, or static, depending on whether they require running the test. Rerunning the suite is practiced at scale at Google and Microsoft \cite{lam2020study, micco2017state}, but its cost grows with the rerun budget, and even five reruns reveal only about 88\% of flaky tests \cite{lam2020understanding}. \Den{Conferir este número na fonte (Lam et al., ISSRE 2020).} iDFlakies reruns suites in different orders to separate order-dependent tests from the others \cite{lam2019idflakies}. DeFlaker flags tests that fail without exercising modified code \cite{bell2018deflaker}; Shaker injects stress to expose concurrency-related flakiness \cite{silva2020shake}. Several tools target .NET directly. RootFinder analyzes logs of passing and failing runs to root-cause flaky tests at Microsoft \cite{lam2019root}, TSVD detects thread-safety violations in C\# code during testing \cite{li2019tsvd}, and FlakeRepro \cite{leesatapornwongsa2022flakerepro} combines static and dynamic analysis to reproduce the failure of a concurrency-related flaky test: across 31 flaky tests from 22 Microsoft projects, it reproduced 26. FlakeSync repairs flaky tests caused by asynchronous waits \cite{rahman2024flakesync}. These tools require the test binaries and, in most cases, an observed failure, so they act during debugging after flakiness has been found, whereas our approach predicts which tests are likely flaky from source code alone and can feed into them.

\myparagraph{Vocabulary-based prediction} Pinto et al.~\cite{pinto2020vocabulary} trained classifiers on the vocabulary of test code (F1 = 0.95 in Java). Camara et al.~\cite{camara2021revisiting} replicated the study with more classifiers and found that recall falls to about 0.02 in an inter-project evaluation; Haben et al.~\cite{haben2021replication} carried it to Python and observed that a time-based split lowers performance; and Soratto and Silva~\cite{soratto2023vocabulary} reported a cross-project F1 of 0.07 for JavaScript. Ahmad et al.~\cite{ahmad2025understanding} compared five languages, and Medeiros et al.~\cite{medeiros2026vocabulary} studied Swift and found a vocabulary of asynchronous constructs and stability markers. In an industrial setting, Berndt et al.~\cite{berndt2023vocabulary} found that the predictions, despite high F-measures, were not actionable for developers. Other static predictors reach similar conclusions: FlakeFlagger \cite{alshammari2021flakeflagger} reports that tokens of the test source were not effective predictors and relies on behavioral features such as coverage and execution time, FLAST \cite{verdecchia2021know} classifies tests by similarity to known flaky ones, and Pontillo et al.~\cite{pontillo2022static} and Gruber et al.~\cite{gruber2023practical} evaluate static and history-based features under more realistic conditions. Barbosa et al.~\cite{barbosa2023test} show that root causes vary by ecosystem. \Den{Verificar se C\# está entre as linguagens estudadas por Barbosa et al.; se estiver, citar os resultados para C\# e ajustar a afirmação de novidade.} Table~\ref{tab:literatura} summarizes the published intra-project results. They are consistently high, whereas cross-project results are consistently low. No vocabulary-based study covers C\#.

\begin{table}[!htbp]
\caption{Published intra-project results for vocabulary-based flaky-test prediction. N: tests in the evaluated dataset.}
\label{tab:literatura}
\centering
\footnotesize
\setlength{\tabcolsep}{3pt}
\begin{tabular}{@{}llrcccl@{}}
\toprule
\textbf{Study} & \textbf{Lang.} & \textbf{N} & \textbf{P} & \textbf{R} & \textbf{F1} & \textbf{Protocol} \\
\midrule
Pinto et al.~\cite{pinto2020vocabulary} & Java & 2,806 & 0.99 & 0.91 & 0.95 & 80/20 \\
Camara et al.~\cite{camara2021revisiting} & Java & 2,806 & 0.98 & 0.89 & 0.94 & 80/20 \\
Haben et al.~\cite{haben2021replication} & Python & 20,641 & 0.94 & 0.74 & 0.80 & 80/20 \\
Soratto and Silva~\cite{soratto2023vocabulary} & JS & 868 & --- & 0.72 & 0.82 & 80/20 \\
Medeiros et al.~\cite{medeiros2026vocabulary} & Swift & 182 & 0.92 & 0.83 & 0.87 & $k$-fold \\
\midrule
This work (Extra Trees) & C\# & 60$^{a}$ & 0.87 & 0.63 & 0.69 & 5-fold $\times$ 10 \\
\bottomrule
\end{tabular}

\vspace{0.2em}
\raggedright\footnotesize $^{a}$30 flaky and 30 stable tests per undersampling draw, out of 12,665 labeled methods.
\end{table}

\myparagraph{Positioning} Cross-project evaluation is not new \cite{camara2021revisiting, soratto2023vocabulary, gultekin2026limits}, and we adopt it rather than claim it. This work differs from previous ones in three respects. First, we document the pitfalls of rerun studies in .NET (Table~\ref{tab:armadilhas}) and use the declaring method, with inheritance resolved, as the unit of analysis instead of the test case. Second, we compare two rerun rounds and measure how much the labels, and the conclusions drawn from them, depend on the round. Third, we evaluate the ranking produced by the model at the natural rate of flaky tests, with average precision and the share of flaky tests found in the top of each repository's ranking, and we separate frequent, general markers from rare, specific ones instead of treating the vocabulary as a single ranking. \Den{Se o paper de Go estiver publicado (ou no arXiv) antes da submissão, citá-lo em terceira pessoa (atenção ao double-blind) e alinhar as conclusões: lá o nome do repositório empata com o vocabulário, e um revisor que conheça os dois vai cobrar o mesmo baseline aqui.}

\myparagraph{Learned features} Pre-trained models such as CodeBERT \cite{fatima2022flakify} and prompting-based use of LLMs \cite{berndt2026can} provide mixed evidence: learned features help more in categorizing root causes than in binary prediction \cite{rahman2024quantizing}, accuracy drops considerably under realistic imbalance \cite{rahman2025understanding}, and code-based detectors lose most of their advantage once the evaluation separates projects \cite{gultekin2026limits}. Vocabulary features remain cheap and interpretable. \Den{Em 2027 os revisores vão pedir ao menos um baseline de modelo pré-treinado (Flakify/CodeBERT) ou de LLM com prompt (como Berndt et al. 2026 e a RQ4 do paper de Go). Com 30 positivos o fine-tuning é inviável, mas um LLM zero-shot sobre os 30 flaky e uma amostra de estáveis é barato e responde à pergunta ``por que não LLM?''.}
